\documentclass[10pt,letterpaper]{article}

\usepackage[top=0.4in,bottom=0.4in,left=0.5in,right=0.5in]{geometry}
\usepackage{enumitem}
\usepackage{titlesec}
\usepackage[hidelinks]{hyperref}
\usepackage{textcomp}

\pagestyle{empty}
\setlength{\parindent}{0pt}
\linespread{0.95}

% Section headings: bold, small-caps-ish, with a rule underneath
\titleformat{\section}{\normalsize\bfseries}{}{0pt}{}[\vspace{-0.6em}\rule{\textwidth}{0.6pt}]
\titlespacing{\section}{0pt}{0.4em}{0.15em}

% Tight bullet lists
\setlist[itemize]{leftmargin=1.1em,itemsep=0pt,parsep=0pt,topsep=0.08em}

% Role heading: title (bold) left, dates right; org left, location right
\newcommand{\role}[4]{%
  \textbf{#1}\hfill #2\\
  \textit{#3}\hfill \textit{#4}\\[0.1em]
}

\begin{document}

%% ─── Header ───────────────────────────────────────────────────────────────
\begin{center}
  {\LARGE\bfseries Edward Lin}\\[0.3em]
  +1 (647) 996-5538 $|$ Toronto, Ontario, Canada $|$ Currently in Taipei, Taiwan\\[0.25em]
  \href{mailto:edwardlin.dev@outlook.com}{edwardlin.dev@outlook.com} $|$
  \href{https://github.com/Edward9292}{github.com/Edward9292} $|$
  \href{https://www.linkedin.com/in/edwardlindev/}{linkedin.com/in/edwardlindev}
\end{center}

%% ─── Experience ───────────────────────────────────────────────────────────
\section*{Professional Experience}

\role{Software Engineer}{Sept 2026 -- Present}{Google -- Cloud AI, Accelerator Software (ASIC-SW)}{Taipei, Taiwan (Hybrid)}
\begin{itemize}
  \item Develop the C++ golden reference model (GRM) for Google's Tensor Processing Unit (TPU) AI
        accelerators, a functional software simulator of the hardware used as the reference for
        validating silicon behavior.
  \item Work across the accelerator software toolchain supporting TPU architecture and verification teams.
\end{itemize}

\vspace{0.25em}
\role{Software Engineer}{Nov 2023 -- Aug 2026}{DAS Inc. -- Developer of DAS Trader}{Markham, Ontario, Canada}
\begin{itemize}
  \item Independently designed and built a GPU-accelerated rendering engine from scratch using C++, OpenGL,
        and GLSL, replacing a 20-year-old CPU-based system, achieving a 15$\times$ increase in rendering
        speed, 42\% reduction in CPU utilization, and enabling real-time visualization of high-frequency
        trading data.
  \item Leveraged RenderDoc for frame capture, shader debugging, and GPU performance analysis, identifying
        bottlenecks and optimizing render passes for maximum efficiency.
  \item Developed multi-GPU support with efficient context management, ensuring low-latency rendering across
        multiple GPUs and monitors for enhanced scalability.
  \item Introduced multi-threaded and parallel rendering from the ground up, enabling simultaneous real-time
        rendering of multiple stock charts with consistent frame rates under heavy data loads.
  \item Optimized real-time stock data processing by designing preprocessing logic from scratch, achieving
        130K+ updates per second and reducing end-to-end latency by 26\%.
  \item Collaborated directly with users to translate real-world workflow feedback into visualization and
        usability improvements.
\end{itemize}

\vspace{0.25em}
\role{Software Engineer Intern}{May 2021 -- Apr 2022}{Oracle Corporation}{Toronto, Ontario, Canada}
\begin{itemize}
  \item Contributed to the backend development of Oracle Responsys, working closely with the team on a
        large-scale codebase that powers a leading cross-channel marketing automation platform.
  \item Developed and optimized Java-based microservices using the Spring Framework, collaborating on 6+ new
        features that improved the efficiency of filtering millions of emails and SMS messages.
  \item Improved the data filtering algorithm and restructured data handling in existing code, reducing
        processing bottlenecks and achieving a 23\% speed increase for large dataset operations.
  \item Enhanced backend server reliability by increasing unit test coverage by 15\%, using JUnit to improve
        code maintainability and minimize regression issues.
  \item Resolved 12+ high-priority bugs that caused filtering algorithms to perform unexpectedly and slow
        with large datasets, improving backend reliability and optimizing system performance.
\end{itemize}

%% ─── Skills ───────────────────────────────────────────────────────────────
\section*{Skills}
\textbf{Programming Languages:} C++, C, Java, Python, JavaScript, R, Shell Scripting\\
\textbf{Frameworks \& Libraries:} React Native, React.js, Node.js, Django, PyTorch, TensorFlow, OpenCV, OpenMP\\
\textbf{Graphics \& GPU Computing:} OpenGL, GLSL, RenderDoc\\
\textbf{Tools \& Databases:} Git, SQL, JUnit

%% ─── Education ────────────────────────────────────────────────────────────
\section*{Education}
\textbf{University of Toronto}\hfill Toronto, Ontario, Canada\\
\textit{Honours Bachelor of Science in Computer Science with a Statistical Minor}\hfill \textit{Sept 2018 -- Nov 2023}
\begin{itemize}
  \item Associate Programs: Focus in Artificial Intelligence, Focus in Web and Internet Technologies
\end{itemize}

%% ─── Projects ─────────────────────────────────────────────────────────────
\section*{Projects}
\textbf{Ontario Parks $|$ Mobile Application}\hfill
\href{https://github.com/Edward9292/ontario-parks-app}{github.com/Edward9292/ontario-parks-app}
\begin{itemize}
  \item Collaborated with Ontario Parks to design and develop a mobile audio-tour app using React Native and Node.js.
  \item Implemented user authentication, cross-device data synchronization, and Google Maps API integration
        with real-time location tracking for context-aware audio content delivery.
\end{itemize}

\vspace{0.2em}
\textbf{CartoonGAN $|$ Machine Learning}\hfill
\href{https://github.com/Edward9292/CartoonGAN/blob/main/report.pdf}{github.com/Edward9292/CartoonGAN}
\begin{itemize}
  \item Developed a Generative Adversarial Network (GAN) using Python, PyTorch, OpenCV, and NumPy to
        transform real images into cartoon-style visuals.
  \item Fine-tuned the GAN architecture and hyperparameters, reducing training time by 18\% while
        maintaining high visual quality in generated images.
\end{itemize}

\end{document}
